\renewcommand{\proofname}{Demostración}
\renewcommand{\appendixname}{APÉNDICE}
\chapter*{APÉNDICES}
\appendix
\markboth{\thechapter APÉNDICES}{\thechapter APÉNDICES}
\addcontentsline{toc}{chapter}{APÉNDICES}
\renewcommand{\thesection}{\Alph{section}}
\renewcommand{\theequation}{A.\arabic{equation}}

\section{Resolución ecuaciones de regulación del Lema \ref{lem2:regulation_equation} y Teorema \ref{teo1:OutputRegulation}}
\lhead[\thepage]{\thesection. Resolución ecuaciones de regulación del Lema \ref{lem2:regulation_equation} y Teorema \ref{teo1:OutputRegulation}} \label{appen1}

Se asume el sistema lineal sin distorsión $d(t)=0$ dado por:
\begin{align}
    \dot{x}(t) & = A x(t) + B u(t), \notag \\
    y(t) & = C x(t) + D u(t) 
\end{align}

Se posee el sistema exógeno compuesto únicamente por la referencia dado por:
\begin{align}
    \dot{v}(t) & = A_1 v(t), \notag \\
    r(t) & = C_1 v(t)  
\end{align}

Una vez descrito el sistema lineal sin distorsión y el sistema exógeno de la referencia ($E=0$ y $F=-C_1$ respectivamente de Teorema \ref{teo1:OutputRegulation}), las ecuaciones de regulación que caracterizan al problema de Output Regulation se definen como:
\begin{align}
    X A_1 = A X + B U,  \\
    0 = C X + D U - C_1
\end{align}

Mediante el operador $vec(\cdot)$ junto a su propiedad $vec(ABC)=(C^\top \otimes A) vec(B)$ donde $\otimes$ es el producto de Kronecker, se desarrollan las ecuaciones de regulación para encontrar $X$ y $U$:

\begin{itemize}
    \item $ X A_1 = A X + B U $
    
    \begin{align} \label{eq:apendixA_regulation1}
        vec(A X) + vec(B U) & = vec(X A_1) \notag \\
        vec(A X I_1) - vec(I_2 X A_1) + vec(B U I_3) & = 0 \notag \\
        (I_1\otimes A) vec(X) - (A_1^\top \otimes I_2) vec(X) + (I_3 \otimes B) vec(U) & = 0 \notag \\
        \left[(I_1\otimes A) - (A_1^\top \otimes I_2)\right] vec(X) + (I_3 \otimes B) vec(U) & = 0
    \end{align}
    con $I_1$, $I_2$ y $I_3$ matrices identidad de dimensiones apropiadas.

    \newpage
    \item $0 = C X + D U - C_1$

    \begin{align} \label{eq:apendixA_regulation2}
        vec(CX) + vec(DU) - vec(C_1) & = 0 \notag \\
        vec(C X I_1) + vec(D U I_3) & = vec(C_1) \notag \\
        (I_1 \otimes C) vec(X) + (I_3 \otimes D) vec(U) & = vec(C_1)
    \end{align}
\end{itemize}

Juntando las ecuaciones \ref{eq:apendixA_regulation1} y \ref{eq:apendixA_regulation2} en forma matricial tenemos:

\begin{equation}
    \underbrace{\begin{bmatrix}
        \left[(I_1\otimes A) - (A_1^\top \otimes I_2)\right] & (I_3 \otimes B) \\ (I_1 \otimes C) & (I_3 \otimes D)
    \end{bmatrix}}_{M}
    \underbrace{\begin{bmatrix}
        \bar{X} \\ \bar{U}    
    \end{bmatrix}}_{x} = 
    \underbrace{\begin{bmatrix}
        0 \\ \bar{C_1}
    \end{bmatrix}}_{y}
\end{equation}

donde $\bar{X}=vec(X)$, $\bar{U}=vec(U)$  y $\bar{C_1}=vec(C_1)$.

Como se puede apreciar, esta ecuación matricial $(M x = y)$ se puede resolver para $x$ mediante división matricial de la forma:
\begin{equation}
    x = M^{-1} y
\end{equation}
\begin{remark}
    El comando para resolver la ecuación matricial en Matlab es $M \backslash y$.
\end{remark}

Finalmente con $x$ obtenido, se utiliza el operador inverso $vec(\cdot)^{-1}$ para obtener las soluciones originales $X$ y $U$ buscadas:

\begin{equation}
    vec^{-1}(x) = vec^{-1}\begin{bmatrix} \bar{X} \\ \bar{U} \end{bmatrix} = \begin{bmatrix}
        X \\ U
    \end{bmatrix}
\end{equation}





